\section{PRINCIPIOS BÁSICOS DE PROTECCIÓN CONTRA RADIACIÓN EXTERNA}

\subsection{Distancia y Tiempo Como Medios de Protección (2 hr)}

El alumno calculará la rapidez de exposición a diferentes distancias de una fuente de actividad conocida. Calculará la exposición en términos de rapidez de exposición y el tiempo. Calculará el tiempo que puede permanecer una persona en un campo de radiación conocido, para recibir una exposición o una dosis predeterminada a distintas distancias de la fuente. Deberá quedar suficientemente familiarizado con este tipo de cálculos como para hacerlos sin dificultad en forma rutinaria.

\subsubsection{PRÁCTICA NO.~6.- Aplicación de la ley del cuadrado inverso. (1:30 hr)}

\subsection{El Blindaje Como Medio de Protección (1 hr)}

El alumno podrá determinar la atenuación en un haz de radiación gamma al atravesar un blindaje de espesor y material conocidos empleando para ello una gráfica semilogarítmica y aplicando el concepto de una capa hemirreductora. Calculará y diseñará blindajes para determinar espesores de blindaje suficientes, a fin de producir la atenuación deseada.

\subsubsection{PRÁCTICA NO.~7.- Atenuación de la radiación gamma, determinación de la capa hemirreductora utilizando diversos materiales (plomo, concreto, ladrillo, agua). (1 hr)}

\subsection{Protección Contra la Contaminación Externa y la Ingestión o la Inhalación. (0:30 hr)}

El alumno reafirmará sus conceptos de contaminación y podrá explicar la deferencia entre contaminación externa, ingestión e inhalación, así como la forma en que pueden producirse y los medios para evitarla.

\subsection{La Protección Radiológica Durante el Trabajo. (1 hr)}

Identificará las situaciones de mayor riesgo durante el trabajo radiográfico. Describirá secuencialmente las medidas de protección que debe tomarse en una jornada de trabajo.